\section{从圆锥曲线反演得到帕斯卡蜗线}
\markboth{70. 从圆锥曲线反演得到帕斯卡蜗线}{70. 从圆锥曲线反演得到帕斯卡蜗线}

以焦点为极点的圆锥曲线极坐标方程可以写成
\[
  \rho = \frac{p}{1+\varepsilon\cos\theta},
\]
其中当离心率 $\varepsilon>1$ 时为双曲线，$\varepsilon=1$ 时为抛物线，$\varepsilon<1$ 时为椭圆（特例当 $\varepsilon=0$ 时为圆）。以极点为反演极，任意正数 $R^2$ 为反演幂，点 $(\rho,\theta)$ 变成 $\left(\frac{R^2}{\rho},\theta\right)$，因而所得曲线的方程是
\[
  \rho = \frac{R^2}{p}(1+\varepsilon\cos\theta).
\]
记
\[
  \frac{R^2}{p} = l,\quad \varepsilon l = 2r,
\]
方程化为
\[
  \rho = 2r\cos\theta + l.
\]
我们在第 30 节中已经见到过上述方程，由它表示的曲线是帕斯卡蜗线：当 $2r>l$（$\varepsilon>1$）时曲线带有绕扣，$2r=l$（$\varepsilon=1$）时曲线有尖点（心脏线），$2r<l$（$\varepsilon<1$）时曲线没有绕扣和尖点。这样就证明了，以圆锥曲线的焦点（即极点）为反演极，作反演变换，结果把这条圆锥曲线变成帕斯卡蜗线。三种不同类型的圆锥曲线（双曲线、抛物线、椭圆），分别对应于三种不同类型的帕斯卡蜗线。
